\section{Validation Methodology}
\label{sec:valmethod}

\subsection{Why not sensitivity/specificity}
Computing true per-case sensitivity/specificity requires a radiologist's per-case measurements as ground
truth. No public dataset provides this for cervical MRI (Section~\ref{sec:data}). We therefore define
``validated'' operationally:
\begin{quote}
\emph{Healthy cases sit on the normal side of each cited threshold and never cross; a pathology-enriched
cohort shifts/crosses into the abnormal side --- with documented limits.}
\end{quote}
This is a \textbf{threshold-crossing $+$ distribution-separation} design. It demonstrates that each
measurement reads the correct range, without claiming per-case diagnostic accuracy.

\subsection{The load-bearing idea: disease-agnostic geometry}
Each measurement reports a physical quantity (mm, degrees, ratio, present/absent) whose normal range is
anchored on healthy-population literature, independent of disease. Three consequences make a single
demonstration cohort sufficient for the groups it exercises:
\begin{enumerate}[leftmargin=*]
  \item Nothing is \emph{trained} on the unhealthy cohort (segmenters are frozen, thresholds are cited)
  --- there is no overfitting-to-spondylosis risk; the cohort is a demonstration set, not a training set.
  \item A measurement crosses its threshold whenever the \emph{anatomy} is abnormal, regardless of cause,
  so abnormal-detection generalizes across diseases producing the same anatomical change. We do not need
  a separate validation dataset per disease to trust that the canal-AP measurement trips on a narrow canal.
  \item The healthy anchor is the real guarantee: if the normal range is calibrated, a value outside it is
  a flagged finding (for review, not diagnosis).
\end{enumerate}
The corollary defines the genuine gaps: findings the system does not \emph{measure} at all (e.g.\ a
tumour mass that does not change canal/cord/disc/vertebra geometry; foraminal stenosis on sagittal-only
MRI) would be missed; and an abnormal value is not by itself a symptomatic disease (specificity overlap),
hence flag-for-review.

\subsection{Cohorts, statistics, and figures}
Healthy $=$ 12 Spine-Generic controls; unhealthy $=$ 10 MMCSD symptomatic cases (5 CSM, 5 CSR). For each
case we compute a per-case summary metric per measurement (e.g.\ the across-level minimum canal AP), and
compare the healthy and unhealthy distributions with a two-sided Mann--Whitney $U$ test (a
non-parametric test appropriate for small $n$ and non-normal data). Per-measure strip plots (healthy vs.\
unhealthy, with the median and the $p$-value) are generated with matplotlib (open source). All figures and
the numeric results are reproducible from \texttt{run\_validation\_master.py}.

\subsection{What a correct result looks like per group}
By construction we expect: strong separation where the disease lives (canal/cord stenosis $\rightarrow$
Group 3); weaker or absent separation where the metric is less specific (alignment $\rightarrow$ Group 4);
a null result where the cohort lacks the pathology (vertebral compression $\rightarrow$ Group 1 on a
non-compression cohort); and the harness should catch any metric that points the wrong way (a built-in
bug detector). Section~\ref{sec:results} shows each of these occurred --- and, importantly, that the
alignment lead which looked directional at small $n$ \emph{dissolved to a non-discriminator} once a
representative, balanced cohort was used.
